\documentclass[11pt,a4paper]{article}

\usepackage[british]{babel}
\usepackage[utf8]{inputenc}

\usepackage{amsmath}
\usepackage[
    % n,
    % advantage,
    operators,
    sets,
    % adversary,
    % landau,
    probability,
    % notions,
    % logic,
    % ff,
    % mm,
    % primitives,
    % events,
    % complexity,
    % oracles,
    % asymptotics,
    % keys
]{cryptocode}
\usepackage{graphicx}
\usepackage{hyperref}
\usepackage{pdflscape}
\usepackage{tabularx}
\usepackage{tikz}

\createprocedureblock{procb}{center,boxed}{}{}{}
\usetikzlibrary{positioning}

\newcommand{\FMDRegister}{\mathsf{FMD.Register}}
\newcommand{\FMDLogin}{\mathsf{FMD.Login}}
\newcommand{\FMDUploadData}{\mathsf{FMD.UploadData}}
\newcommand{\FMDDownloadData}{\mathsf{FMD.DownloadData}}

\newcommand{\AESGCMKeyGen}{\mathsf{AES\text{-}GCM.Keygen}}
\newcommand{\AESGCMEnc}{\mathsf{AES\text{-}GCM.Enc}}
\newcommand{\AESGCMDec}{\mathsf{AES\text{-}GCM.Dec}}
\newcommand{\Argon}{\mathsf{Argon2id}}
\newcommand{\HKDF}{\mathsf{HKDF}}
\newcommand{\SHATwoFiveSix}{\mathsf{SHA2\text{-}256}}
\newcommand{\SHAFiveTwelve}{\mathsf{SHA2\text{-}512}}

\usepackage{titling}
\setlength{\droptitle}{-5em} % reduce space above title

\title{The FMD Server Protocol \\[1em] \includegraphics[scale=0.1]{logo.png}}
\author{Thore Göbel}
\date{Version 2.0 -- 2026-09-05}


\begin{document}

\maketitle

\tableofcontents


\newpage

\section{Overview}

FMD (``Find My Device'', \href{https://fmd-foss.org/}{fmd-foss.org}) is a libre
and open source system for remotely finding and controlling one's device.
\\

Users send \emph{commands} to their devices.
For example, to locate the device by requesting its GPS coordinates (\texttt{fmd locate gps}),
to let it ring (\texttt{fmd ring}), to lock it (\texttt{fmd lock <msg>}),
or even to factory-reset it (\texttt{fmd delete <delete-pwd>}).
Commands are strings and use a CLI-like format.

Commands are sent over \emph{transport channels} (or \emph{transports} for short)
to the FMD app running on their Android device.
One basic transport channel is SMS.
Users allow-list the phone numbers that may send commands.
\\

A more advanced transport channel is \emph{FMD Server}.
Users register their devices with the server, and then use the web interface
to send commands to their device.
The web interface also shows the responses, such as a map with the locations
returned by the device, and pictures taken by the device.

Fundamentally, FMD Server is just a backend with a database and a REST API.
Notably, it deals with sensitive data (device locations, pictures of the device's surroundings)
and it exposes sensitive functionality (remotely controlling people's devices).
Therefore, data protection and access control are crucial.

FMD Server is self-hostable, which is the recommended option.
Self-hosted instances may still be shared by multiple users
(for example, a family admin running it for multiple people).
The FMD project provides a hosted instance for users who don't have the means
to self-host FMD Server. \footnote{\url{https://server.fmd-foss.org/}}
In both cases, a natural goal for FMD Server is to prevent a rogue server administrator
from controlling the devices of the users registered on the server.
\\

This whitepaper describes the threat model for FMD Server, its security goals,
and the cryptographic protocol designed to achieve these goals.
Together, they form the \emph{FMD Server Protocol}.


\newpage

\section{Threat model} \label{sec:threat-model}

This section describes the threat model for the FMD Server Protocol.
That is, how much power we grant the adversary.

\paragraph{Assumption 1: Malicious backend}
We assume that the host that runs the backend code is compromised.
That is, the adversary can read and write the database,
read and write all incoming and outgoing network requests, and change the backend code.

This is motivated by shared hosting environments:
users should be able to use an FMD Server instance hosted by people they don't trust.

\paragraph{Assumption 2: Network adversary}
Since the adversary has access to the server, they also have all properties of a Dolev-Yao adversary.
That is, they can read, modify, reorder, replay, and drop messages sent over the network.

\paragraph{Assumption 3: Secure client}
We assume that the adversary cannot compromise the client code.
That is, the client adheres to the protocol and does not reveal secrets to the adversary.
\\

For the Android app, this can be ensured by means of inspecting the source code,
verifying signed build artifacts, and Reproducible Builds.

For the web client, this is a huge, not currently realistic simplification.
This is because the frontend web client is served by the same server that hosts the backend code.
And that server we assumed above to be compromised!
Unfortunately, this problem is inherent to the web as a platform.
Therefore, we will set this problem aside.
For example, one workaround would be to ship the web client as an Electron-based desktop app.
\\

Given FMD's use case of lost or stolen devices,
the adversary may have physical access to the device.
The secure client assumption means that the adversary cannot physically tamper
with the device in a way that compromises the FMD client application.


\newpage

\section{Security goals}

This section describes the security goals that the FMD Server Protocol (described
in \autoref{sec:protocol}) aims to achieve.

\paragraph{Goal 1: Confidentiality}
The adversary does not learn sensitive data generated by an honest client,
such as the content of locations and pictures.

The adversary does not learn the arguments of the commands that are executed.
For example, the adversary does not learn the content of the lockscreen message
in the command \texttt{fmd lock "Please call me at +1 23 456"}.

\paragraph{Goal 2: Integrity}
All data that an honest client accepts has been previously generated by a honest client of the same account.
The adversary cannot inject new data or swap data between accounts.
This includes locations/pictures/etc. as well as commands.

\paragraph{Goal 3: No replay}
The adversary cannot replay a command to cause a client to execute a command more than once.


\subsection{Non-goals}\label{sec:non-goals}

The FMD Server Protocol is \emph{not} designed to protect against the scenarios listed below.
These are explicitly out-of-scope.

\paragraph{Non-goal 1: Command Type Confidentiality}
The adversary is allowed to learn what types of command are executed (locate, picture, ...).
This is necessary for commands that return data.
The server needs to know the data type so that the web client can query a list of locations
(as opposed to a list of pictures).

\paragraph{Non-goal 2: Availability}
We exclude availability and denial-of-service attacks.
While service unavailability is extremely frustrating when users lose their device and need the service,
a malicious server operator or on-path network adversary can always and trivially block users.
Solving this would require peer-to-peer solutions; yet the design of FMD Server
is intentionally ``centralised'' on a single, self-hostable instance.

\paragraph{Non-goal 3: Data completeness}
The server is allowed to remove data (locations, pictures, commands).
A malicious server may use this to selectively hide data from a client.
For example, it may only show a subset of locations or it may not deliver a command.

The motivation for this is to enable the server to autonomously prune data
(for example, to keep only the last 10 pictures),
and to keep the protocol simple.
We leave it to future work to include integrity protection for this in the protocol.




\newpage

\section{Protocol}\label{sec:protocol}

This section describes the FMD Server Protocol.

\subsection{Notation}

In the protocol descriptions that follow, the following notations are used:

\begin{itemize}
    \item[$ || $] \quad bitstring concatenation
    \item[$ \bin^{l} $] \quad set of bitstrings of length $l$
    \item[$ v \sample S $] \quad value $v$ is sampled uniformly at random from set $S$
    \item[$ v \gets e $] \quad value $v$ is assigned the result of expression $e$
\end{itemize}

\subsection{Participants}

There are two protocol participants: the \textbf{client} $A$ (Android, web, etc.)
and the \textbf{server} $S$ (backend).

Where necessary, we distinguish between the \emph{device client}
(that is \emph{receiving} commands and is uploading locations/pictures/etc)
and the \emph{web client} (that is \emph{sending} commands and viewing uploaded data).

\subsection{Fundamental concepts}

First, we describe the fundamental concepts that are necessary to understand the design of the protocol.

\paragraph{Protected data}
There are different types of client-to-client data that require protection:

\begin{itemize}
    \item \emph{Commands:}
        Uploaded by the web client to the FMD Server backend.
        Retrieved by the device client (FMD Android) and executed.
    \item \emph{Locations:}
        Uploaded by the device client to FMD Server.
        Retrieved and displayed by the web client.
    \item \emph{Pictures:} Same as locations.
    \item \emph{Free-text responses:}
        \footnote{This is not yet implemented at the time of writing, but is documented here as a future feature.}
        Uploaded by the device client to FMD Server
        in response to running a command.
        This can be a free-text result string or an error message.
        This is the catch-all option for commands other than \texttt{locate} and \texttt{camera}.
\end{itemize}

In the future, there may be other device clients (desktop, CLI), but the ideas are the same.
There may also be new types of structured data, but the ideas (key hierarchy, encryption algorithm)
from locations/pictures extend to these other data types.

\paragraph{Unprotected data}
The following data is out-of-scope for the protocol:

First, server-to-client data that the server generates and sends to the client.
For example, when the server detects multiple failed login attempts
it sends a warning to any logged-in client.
Second, client-to-server data.
For example, the client uploads a push URL, which the server uses to wake up the device
using UnifiedPush\footnote{\url{https://unifiedpush.org/}}.

This data is \emph{not} protected by the FMD Server Protocol;
it solely relies on the protections provided by TLS.
Only client-to-client data is covered by the protocol.

\paragraph{Accounts}
Each \emph{account} on FMD Server has a username and passphrase.
The Android device client is used to register a new account.
The web client can be used to log into an account.
Every account has a \emph{username} $U$.
\\

There should only be a single device client for every account.
Multi-device accounts are future work (see \autoref{sec:future-work}).
\footnote{
The Android client can be used to log into an existing account.
It is the responsibility of the user to make sure that they don't log into the
same account on multiple devices.

If multiple Android devices are logged into the same account, the behaviour is undefined.
For example, the push URL of all but the latest device will be dropped.
Also, the user will not be able to distinguish from which device data is coming from.
}

\paragraph{Passphrase}
The passphrase $P$ acts as the root of trust.
The protections provided by the security protocol hinge on this passphrase.
\textbf{If the adversary learns the passphrase, the protocol makes no security guarantees.}

Therefore, the user needs to strike a balance between a passphrase
that is strong enough against a brute-forcing adversary, and a passphrase that they can
recover in the stressful situation of having lost their phone.
For example, they may remember it, look it up in the password manager on their laptop,
or have it on a piece of paper under their pillow.

The motivation for using a passphrase is to allow for flexibility when the user loses their device.
In such a situation, they need to be able to ask a friend -- or even a random stranger on the street -- for help.
Users need to be able to access their FMD Server account via the web client from any browser.
This would not always be possible with software passkeys or hardware security keys
(for example, your bag was stolen, and the bag contained both your phone and your keyring with the hardware security key).


\paragraph{User authentication}
The user authenticates to the FMD Server backend to be allowed to make API requests to access an account.

Note that malicious read/write access to an account is allowed by the threat model
(because we assumed the server to be malicious).
However, in practice, we clearly want the server to prevent random people
on the internet from modifying account data.
This includes data protected by the client-to-client protocol described in this paper,
but also data that is client-to-server.
For example, only an authenticated user should be able to change the push URL
that FMD Server uses to wake up a device when a command is pending.


\subsection{Passphrase hashing}\label{sec:passphrase-hashing}

The client uses a password-based key derivation function to derive strong symmetric
secrets from the likely human-chosen, low-entropy passphrase.

\paragraph{Client-side hashing}
Specifically, the client derives an initial \emph{password key} $K_{pwk}$ from the passphrase,
as shown in \autoref{fig:passphrase-hashing}.
It is used as the start of the key schedule (\autoref{sec:key-schedule}).
It never leaves the client.

As the password-based KDF, $\Argon$ is used with the parameters listed in \autoref{tab:argon-client}.
These parameters were chosen using the considerations of section 4 of
\href{https://datatracker.ietf.org/doc/html/rfc9106}{RFC 9106}.
$m$ was chosen to account for the fact that the clients are mobile devices.
In our testing, these values of $m$ still ran in reasonable time (under two seconds),
whereas higher powers of 2 took significantly longer or caused out-of-memory (OOM)
issues on older/lower-end Android devices.

\begin{figure}[h]
    \centering
    \makebox[0cm][]{
        \begin{tikzpicture}[thick]

            \node (P) {``fmd\_v2\_password'' \concat\ $h(U)$ \concat\ passphrase $P$};

            % hashing password key
            \node[draw, dotted] (Argon) [below=of P] {$\Argon$};
            \node (S) [left=of P] {salt $S$};
            \node (PMK) [below=of Argon] {password key $K_{pwk}$};
            \node (Argoninfo) [right=0.3cm of Argon] {};

            \draw[-] (P) -- (Argon);
            \draw[-] (S) -- node {} (Argon);
            \draw[->] (Argon) -- node {} (PMK);

        \end{tikzpicture}
        }

        \caption{Passphrase hashing (client-side)}
        \label{fig:passphrase-hashing} % MUST come after \caption!
    \end{figure}


\paragraph{Server-side hashing}
From $K_{pwk}$, the client derives a separate \emph{authentication key} $K_{auth}$ with $\HKDF$
(see the key schedule in \autoref{sec:key-schedule}).
The server receives $K_{auth}$ from the client.
It treats this key as a password-like symmetric authentication token.
The server hashes it again with $\SHAFiveTwelve$
and compares the resulting hash $h(K_{auth})$ against the hash stored in the database.
If they match, the user is authenticated.

\begin{table}[]
    \centering
    \begin{tabular}{|l|l|l|}
    \hline
    Parameter & Value  & Description \\ \hline
    $t$       & 1      & iterations\\
    $p$       & 4      & lanes \\
    $m$       & $2^{17}$ KiB = $128$ MiB & memory size \\
              & 256 bit = 32 byte & hash length \\
              & 128 bit = 16 byte & salt length \\ \hline
    \end{tabular}

    \caption{Argon2id parameters for client-side hashing}
    \label{tab:argon-client} % MUST come after \caption!
\end{table}



\subsection{Key Schedule}\label{sec:key-schedule}

\autoref{fig:key-schedule} and \autoref{fig:key-schedule-data} show
the key schedule of how keys are derived.
These figures are intended to be read together with the algorithms using these keys,
described in \autoref{sec:subprotocols}.

The key schedule is split into two halves.
The top half derives an \emph{authentication key} $K_{auth}$
and a \emph{pre-master key} $K_{pmk}$ from the user's passphrase.
These change whenever the user decides to rotate their passphrase.

The bottom half derives child keys from a \emph{master key} $K_{master}$.
The master key is generated once during account creation, and afterwards never changes.

The child keys act as key encryption keys (KEKs).
There is one child key per data type: a location key $K_{loc}$, a picture key $K_{pic}$,
a command key $K_{cmd}$, and so on.

The individual data items are encrypted using unique, per-item data encryption keys (DEKs).
These are freshly generated for every item, and are encrypted under their respective KEK.
This is shown in \autoref{fig:key-schedule-data}.

\begin{landscape}
    % XXX: This also shifts the page number
    \addtolength{\hoffset}{-2.5cm}
    \addtolength{\voffset}{+1.75cm}

    \begin{figure}
        \centering
        \small

        \begin{center}
        \begin{tikzpicture}[thick]
            % passphrase to pre-master key
            \node (P) {``fmd\_v2\_password'' \concat\ $h(U)$ \concat\ passphrase $P$};
            \node (S) [left=of P] {salt $S$};
            \node[draw, dotted] (Argon) [below=of P] {$\Argon$};
            \node (PWK) [below=of Argon] {$K_{pwk}$};
            \node (Argoninfo) [right=0.3cm of Argon] {};
            \draw[-] (P) -- (Argon);
            \draw[-] (S) -- node {} (Argon);
            \draw[->] (Argon) -- node {} (PWK);

            % authentication key
            \node[draw, dotted] (KDFauk) [below left=2cm of PWK] {$\HKDF$};
            \node (KDFaukinfo) [below left=2cm of PWK] [right=0.3cm of KDFauk] {salt=$\emptyset$, info=``fmd\_v2\_auth'' \concat\ $h(U)$};
            \node (AUK) [below=2cm of KDFauk] {$K_{auth}$};
            \draw[-] (PWK) -- (KDFauk);
            \draw[->] (KDFauk) -- (AUK);

            % pre-master key
            \node[draw, dotted] (KDFpmk) [right=2cm of KDFaukinfo] {$\HKDF$};
            \node (KDFpmkinfo) [right=0.3cm of KDFpmk] {salt=$\emptyset$, info=``fmd\_v2\_premaster'' \concat\ $h(U)$};
            \node (PMK) [below=2cm of KDFpmk] {$K_{pmk}$};
            \draw[-] (PWK) -- (KDFpmk);
            \draw[->] (KDFpmk) -- (PMK);

            % decrypting the master key
            \node (IVmk) [right=of PMK] {$IV_{master}$};
            \node (EMK) [right=of IVmk] {$C_{master}$};
            \node[draw, dotted] (EMKdec) [below=of IVmk] {$\AESGCMDec$};
            \node (EMKad) [right=0.3cm of EMKdec] {$AD=$ ``fmd\_v2\_master'' \concat\ $h(U)$};
            \node (MK) [below=of EMKdec] {$K_{master}$};
            \draw[-] (IVmk) -- (EMKdec);
            \draw[-] (EMK) -- (EMKdec);
            \draw[-] (PMK) -- (EMKdec);
            \draw[->] (EMKdec) -- (MK);

            % picture keys
            \node[draw, dotted] (KDFpic) [below=2cm of MK] {$\HKDF$};
            \node (KDFpicinfo) [right=0.3cm of KDFpic] {salt=$\emptyset$, info=``fmd\_v2\_kek\_picture'' \concat\ $h(U)$};
            \node (PICK) [below=of KDFpic] {$K_{pic}$};
            \draw[-] (MK) -- (KDFpic);
            \draw[->] (KDFpic) -- (PICK);

            % location keys
            \node (KDFlocinfo) [left=of KDFpic] {salt=$\emptyset$, info=``fmd\_v2\_kek\_location'' \concat\ $h(U)$};
            \node[draw, dotted] (KDFloc) [left=of KDFlocinfo] {$\HKDF$};
            \node (Kloc) [below=of KDFloc] {$K_{loc}$};
            \draw[-] (MK) -- (KDFloc);
            \draw[->] (KDFloc) -- (Kloc);

            % command keys
            \node (KDFcmds) [right=of KDFpicinfo] {$\dots$};
            \draw[-] (MK) -- (KDFcmds);

        \end{tikzpicture}
    \end{center}

    \caption{Key Schedule (part 1)}
    \label{fig:key-schedule} % MUST come after \caption!
\end{figure}



\begin{figure}
    \centering

    \begin{center}
        \begin{tikzpicture}[thick]

            % location keys
            \node[draw, dotted] (AESKeyGen) {$\AESGCMKeyGen$};
            \node (Kloci) [below=of AESKeyGen] {$K_{k_{id}}$};
            \draw[->] (AESKeyGen) -- (Kloci);

            % encrypt data key
            \node[draw, dotted] (AESEncKey) [below left=4cm of Kloci] {$\AESGCMEnc$};
            \node (IVkey) [above=of AESEncKey] {$IV_{k_{id}}$};
            \node (Kloc) [left=0.5cm of IVkey] {$K_{loc}$};
            \node (ADkey) [right=0.2cm of AESEncKey] {$AD = $ ``fmd\_v2\_dek\_location'' \concat\ $h(U)$ \concat\ $id$ \concat\ $t$};
            \node (EncKey) [below=of AESEncKey] {encrypted location data key $C_{k_{id}}$};
            \draw[-] (Kloc) -- (AESEncKey);
            \draw[-] (IVkey) -- (AESEncKey);
            \draw[-] (Kloci) -- (AESEncKey);
            \draw[->] (AESEncKey) -- (EncKey);

            % encrypt location data
            \node[draw, dotted] (AESEncData) [right=0.1cm of ADkey] {$\AESGCMEnc$};
            \node (IVdata) [above=of AESEncData] {$IV_{d_{id}}$};
            \node (LOCI) [above right=of AESEncData] {$loc_{id}$};
            \node (ADdata) [right=0.1cm of AESEncData] {$AD = $ ``fmd\_v2\_data\_location'' \concat\ $h(U)$ \concat\ $id$ \concat\ $t$};
            \node (EncData) [below=of AESEncData] {encrypted location $C_{d_{id}}$};
            \draw[-] (Kloci) -- (AESEncData);
            \draw[-] (IVdata) -- (AESEncData);
            \draw[-] (LOCI) -- (AESEncData);
            \draw[->] (AESEncData) -- (EncData);

        \end{tikzpicture}
    \end{center}

    \caption{Key Schedule (part 2)}
    \label{fig:key-schedule-data} % MUST come after \caption!
\end{figure}


\end{landscape}




\subsection{Subprotocols}\label{sec:subprotocols}

This section describes the subprotocols of the FMD Server Protocol:
$\FMDRegister$, $\FMDLogin$, $\FMDUploadData$, and $\FMDDownloadData$.

\paragraph{Principles}
The design is guided by the following principles:

\begin{itemize}
    \item Keep it simple.
    \item Use modern algorithms.
    \item Hardcode algorithms and their parameters.
        If we need to switch algorithms, we will release a new protocol version.
    \item Use context binding as much as possible to avoid confusions (in KDFs and AEAD Associated Data).
\end{itemize}

\paragraph{Encoding}
Strings are \textbf{UTF-8} encoded.
Integers are \textbf{big endian} encoded.

Note: For simplicity, the encoding step is usually not shown.

\paragraph{HKDF}
All HKDF calls use $\SHATwoFiveSix$ as the hash function.

\paragraph{Context strings}
Context strings use the following format:

$$ \text{``fmd\_v2\_''} \concat\ \kappa\ \concat\ \text{``\_''}\ \concat\ \tau\ \concat\ h(U)\ \concat\ id\ \concat\ t $$

$\kappa$ is the type (\emph{kek}, \emph{dek}, \emph{data}) of content being derived or encrypted.
$\tau$ is the type of user data.
The username $U$ is hashed. $h$ is fixed to $\SHATwoFiveSix$.
This is to normalise the username string to avoid attacks with maliciously chosen usernames.
The $id$ must be a 128-bit bitstring generated freshly for every data item and uniformly at random.
$t$ is a Unix milliseconds timestamp as a 64-bit integer.

\paragraph{Types of data}
For simplicity, the subprotocols are specified using $\tau$ as the parameter for the data type.
At the time of writing, these are ``location'', ``picture'', ``text'', and ``command''.
Where special handling is needed, this is noted.


\subsubsection{Registration}

The $\FMDRegister$ subprotocol (\autoref{alg:register}) is used to register a new account for user $U$ with the server.

The client derives a pre-master key $K_{pmk}$ from the user's passphrase $P$.
The client locally generates a master key $K_{master}$, symmetrically encrypts it under $K_{pmk}$,
and uploads the result to the server.
This allows passphrase rotation: when the user changes their passphrase,
only the $K_{master}$ needs to be re-encrypted under the new $K_{pmk}'$ and re-uploaded.
This avoids re-encrypting all stored data.


\subsubsection{Login}

The $\FMDLogin$ subprotocol (\autoref{alg:login}) is used by the client to log in, create a new session,
and obtain the $K_{master}$.

For the session, the server returns an $accessToken$ used to authorize subsequent requests
that access that user's data.
Since this is basic client-to-server access control,
we omit the $accessToken$ in the other subprotocol descriptions.
In practice, the client includes the $accessToken$ as an HTTP header field in subsequent requests.


\subsubsection{Uploading data}

The $\FMDUploadData$ subprotocol (\autoref{alg:upload-data}) is used by the client to upload new data to the server.

Every datum is encrypted under a freshly generated \emph{data encryption key (DEK)}.
The DEK is encrypted under a long-term \emph{key encryption key (KEK)}.
Every data type has its own KEK.

The plaintext of every data type internally uses its own structure.
For example, the location may have fields for latitude, longitude, altitude, and so on.
This structure is not relevant for the protocol, we generalise this as a data blob $D$.


\subsubsection{Downloading data}

The $\FMDDownloadData$ subprotocol (\autoref{alg:download-data}) is used by the client to download data from the server.

We assume that the client knows which ids $id_i$ exist for any given data type $\tau$.
In practice, the server provides an unauthenticated list of ids.
This is in line with the non-goals described in \autoref{sec:non-goals}.

Once decrypted, the client parses the data value $D$.
For most types (location, picture, text), it simply displays the data to the user.
For commands, there is additional handling:

\paragraph{Command replay protection}
To prevent replay attacks, the client stores the time $t'$.
When a new command is downloaded, the client compares the new time $t$ against the stored $t'$.
If and only if $t > t'$, the client executes the command, otherwise it rejects it.
Because honest clients choose $t$ as a Unix millisecond timestamp, $t$ should be strictly increasing.

If the client loses its storage (e.g., due to a reinstall), $t'$ is reset,
making replay attacks possible.
This risk is accepted, as device clients are assumed to have long-term storage.

\begin{landscape}

\begin{figure}
    \procb{$\FMDRegister$}{
        \textbf{Client} \> \> \textbf{Server} \\
        \pcln \text{Knows: } U, P \> \> \\
        \\
        \pcln S \sample \bin ^s \pccomment{salt} \> \> \\
        % \pccomment{Compute hashes} \\
        \pcln K_{pwk} \gets \Argon (\text{``fmd\_v2\_password''} \concat\ h(U)\ \concat\ P, \; S) \> \> \\
        \\
        \pcln K_{auth} \gets \HKDF (K_{pwk}, \; salt=\emptyset, \; info = \text{``fmd\_v2\_auth''} \concat\  h(U) ) \> \> \\
        \pcln K_{pmk} \gets \HKDF (K_{pwk}, \; salt=\emptyset, \; info = \text{``fmd\_v2\_premaster''} \concat\  h(U) ) \> \> \\
        \\
        \pcln IV_{master} \sample \bin^{96} \> \> \\
        \pcln K_{master} \sample \bin^{256} \> \> \\
        \pcln AD \gets \text{``fmd\_v2\_master''} \concat\ h(U) \> \> \\
        \pcln C_{master} \gets \AESGCMEnc_{K_{pmk}} (IV_{master}, K_{master}, AD) \> \> \\
        %
        \pcln \> \sendmessageright*{U, S, IV_{master}, C_{master}, K_{auth}} \> \\
        \pcln \> \> \pcif U \text{ exists}: \pcabort \\
        \pcln \> \> hK_{auth} \gets \SHAFiveTwelve (K_{auth}) \\
        \pcln \text{Store: } K_{master} \> \> \text{Store: } U, S, IV_{master}, C_{master}, hK_{auth}
        % TODO: server-side hashing
    }
    \caption{$\FMDRegister$ algorithm}
    \label{alg:register} % MUST come after \caption!
\end{figure}

\begin{figure}
    \procb{$\FMDLogin$}{
        \textbf{Client} \> \> \textbf{Server} \\
        \pcln \text{Knows: } U, P \> \> \text{Knows: } U, S, IV_{master}, C_{master}, hK_{auth} \\
        \\
        \pcln \> \sendmessageright*{U} \> \\
        \pcln \> \sendmessageleft*{S} \> \\
        \pcln K_{pwk} \gets \Argon (\text{``fmd\_v2\_password''} \concat\ h(U)\ \concat\ P, \; S) \> \> \\
        \pcln K_{auth} \gets \HKDF (K_{pwk}, \; salt=\emptyset, \; info = \text{``fmd\_v2\_auth''} \concat\  h(U) ) \> \> \\
        \pcln \> \sendmessageright*{K_{auth}} \> \\
        \> \> hK_{auth}' \gets \SHAFiveTwelve (K_{auth}) \\
        \> \> \pcif hK_{auth}' \neq hK_{auth}: \pcabort \\
        \\
        \> \> accessToken \sample \bin^* \\
        \pcln \> \sendmessageleft*{accessToken} \> \pccomment{authenticates subsequent requests} \\
        \pcln \> \sendmessageleft*{IV_{master}, C_{master}} \> \\
        \pcln K_{pmk} \gets \HKDF (K_{pwk}, \; salt=\emptyset, \; info = \text{``fmd\_v2\_premaster''} \concat\  h(U) ) \> \> \\
        \pcln AD \gets \text{``fmd\_v2\_master''} \concat\ h(U) \> \> \\
        \pcln K_{master} \gets \AESGCMDec_{K_{pmk}} (IV_{master}, C_{master}, AD) \> \> \\
        \\
        \pcln \text{Knows: } K_{master}, accessToken \> \>
    }
    \caption{$\FMDLogin$ algorithm}
    \label{alg:login} % MUST come after \caption!
\end{figure}

\begin{figure}
    \procb{$\FMDUploadData$}{
        \textbf{Client} \> \> \textbf{Server} \\
        \pcln \text{Knows: } U, K_{master} \> \> \text{Knows: } U \\
        \\
        \pcln D \gets \text{data value (location, picture, text, command)} \> \> \\
        \pcln \tau \gets \text{data type} \> \> \\
        \pcln id \sample \bin^{128} \> \> \\
        \pcln t \gets \text{current Unix milliseconds} \> \> \\
        \\
        \pcln K_{\tau} \gets \HKDF (K_{master}, \; salt=\emptyset, \; info = \text{``fmd\_v2\_kek\_''} \concat\ \tau\ \concat\ h(U) \; ) \> \> \\
        \pcln K_{k_{id}} \sample \bin^{256} \> \> \\
        \\
        \pcln IV_{k_{id}} \sample \bin^{96} \> \> \\
        \pcln AD_{k_{id}} \gets \text{``fmd\_v2\_dek\_''} \concat\ \tau\ \concat\ h(U)\ \concat\ id\ \concat\ t \> \> \\
        \pcln C_{k_{id}} \gets \AESGCMEnc_{K_{\tau}} (IV_{k_{id}}, K_{k_{id}}, AD_{k_{id}}) \qquad \pccomment{key encryption} \> \> \\
        \\
        \pcln IV_{d_{id}} \sample \bin^{96} \> \> \\
        \pcln AD_{d_{id}} \gets \text{``fmd\_v2\_data\_''} \concat\ \tau\ \concat\ h(U)\ \concat\ id\ \concat\ t \> \> \\
        \pcln C_{d_{id}} \gets \AESGCMEnc_{K_{k_{id}}} (IV_{d_{id}}, D, AD_{d_{id}}) \qquad \pccomment{data encryption} \> \> \\
        \\
        \pcln \> \sendmessageright*{ U, \tau, id, t, IV_{k_{id}}, C_{k_{id}}, IV_{d_{id}}, C_{d_{id}} } \> \\
        \pcln \> \> \text{Stores: } \langle U, \tau, id, t, IV_{k_{id}}, C_{k_{id}}, IV_{d_{id}}, C_{d_{id}} \rangle
        % \pcln \> \> \\
    }
    \caption{$\FMDUploadData$ algorithm}
    \label{alg:upload-data} % MUST come after \caption!
\end{figure}

\begin{figure}
    \procb{$\FMDDownloadData$}{
        \textbf{Client} \> \> \textbf{Server} \\
        \pcln \text{Knows: } U, K_{master}, t' \> \> \text{Knows: } \{ \langle U, \tau, id, t, IV_{k_{id}}, C_{k_{id}}, IV_{d_{id}}, C_{d_{id}} \rangle _i \}_i \\
        \\
        \pcln \tau \gets \text{data type} \> \> \\
        \pcln K_{\tau} \gets \HKDF (K_{master}, \; salt=\emptyset, \; info = \text{``fmd\_v2\_kek\_''} \concat\ \tau\ \concat\ h(U) \; ) \> \> \\
        \pcln \> \sendmessageright*{U, \tau} \> \\
        %
        \pcln \> \> \{id_i\}_i \gets \text{select all $id$ matching $U$ and $\tau$}\\
        \pcln \> \sendmessageleft*{ \{id_i\}_i } \> \pccomment{server provides unauthenticated list of ids} \\
        \pcln id \gets \text{select from } \{id_i\}_i \> \> \\
        \pcln \> \sendmessageright*{U, \tau, id} \> \\
        %
        \pcln \> \sendmessageleft*{ t, IV_{k_{id}}, C_{k_{id}}, IV_{d_{id}}, C_{d_{id}} } \> \\
        \\
        \pcln AD_{k_{id}} \gets \text{``fmd\_v2\_dek\_''} \concat\ \tau\ \concat\ h(U)\ \concat\ id\ \concat\ t \> \> \\
        \pcln K_{k_{id}} \gets \AESGCMDec_{K_{\tau}} (IV_{k_{id}}, C_{k_{id}}, AD_{k_{id}}) \qquad \pccomment{key decryption} \> \> \\
        \\
        \pcln AD_{d_{id}} \gets \text{``fmd\_v2\_data\_''} \concat\ \tau\ \concat\ h(U)\ \concat\ id\ \concat\ t \> \> \\
        \pcln D \gets \AESGCMDec_{K_{k_{id}}} (IV_{d_{id}}, C_{d_{id}}, AD_{d_{id}}) \qquad \pccomment{data decryption} \> \> \\
        \\
        \pcln \pcif \tau = \text{``cmd''}: \> \> \\
        \pcln \t \pcif t \leq t': \pcabort \> \> \\
        \pcln \t \pcelse : t' \gets t \> \> \\
        \pcln \text{Parse and handle } D \> \>
        % \pcln \> \> \\
    }
    \caption{$\FMDDownloadData$ algorithm}
    \label{alg:download-data} % MUST come after \caption!
\end{figure}

\end{landscape}



\newpage

\section{Future work}\label{sec:future-work}

Below, we list ideas for future extensions of FMD Server.
Especially the last two introduce significant complexity to the specification and an implementation,
and are therefore out-of-scope of this protocol version.

\paragraph{Authentication with an aPAKE}
Instead of deriving a separate authentication key $K_{auth}$ to send to the server as
the authentication value (see \autoref{sec:passphrase-hashing}),
we might instead use an \emph{asymmetric/augmented password-authenticated key exchange (aPAKE)}
for authenticating users.
For example, 1Password\footnote{\url{https://agilebits.github.io/security-design/srp.html}}
and Proton\footnote{\url{https://proton.me/blog/encrypted-email-authentication}} use SRP,
and WhatsApp\footnote{\url{https://www.nccgroup.com/research-blog/public-report-whatsapp-opaque-ke-cryptographic-implementation-review/}}
uses OPAQUE.

\paragraph{Multi-device}
The current protocol is designed with ``one device = one account'' in mind.
A user who has multiple devices (phone, table, laptop) thus needs to register one account for each device.
With multi-device, the user would register a single account to which they can connect multiple devices.

\paragraph{Sharing data between accounts}
Services comparable to FMD -- such as Apple Find My or Google Find Hub --
allow users to share their location with friends and family.
This feature needs asymmetric cryptography (with every account having an identity key pair)
and logic to securely share data selectively between accounts.

At this point, with these sharing capabilities, FMD Server would become
an end-to-end-encrypted cloud storage (even more than it already is).
Encrypted cloud storage is an active area of research\footnote{\url{https://eprint.iacr.org/2024/989}},
with existing schemes often being broken.\footnote{\url{https://eprint.iacr.org/2022/959}, \url{https://eprint.iacr.org/2024/1616}}
Therefore, any such extension requires careful design and formal analysis before deployment.


% \newpage

% \section{Conclusion}

% TODO



\newpage

\appendix

\begin{landscape}

    \section{Version history}

    \begin{tabularx}{\linewidth}{l|l|X|l|l}
        Protocol version & Date published & Description & FMD Server version & REST API version \\
        \hline
        \hline
        0.x & $\approx$ 2021 & Original protocol. Never formally specified. Broken in active adversary model. & 0.x & v1 \\
        \hline
        1.0 & -- & Skipped to align with REST API. & -- & -- \\
        \hline
        2.0 & 2026-09-05 & Initial stable, non-0.x version. & 0.17.0, 1.x & v2 \\
    \end{tabularx}

\end{landscape}

\end{document}
